\cotefile{156-9}
\scanpage{66}
\mspage{62}
% source lines 1--6
que si de plus $\alpha$ est \unsure{surjectif}, i.e. \unsure{bijectif}!

\unsure{Donc} il faut une condition supplémentaire pour pouvoir affirmer la commutation de $\alpha_{\unsure{*}}$ à $C$, ou mieux, la formule $(**)$. On va supposer donné une correspondance en sens inverse
\diagram{\begin{tikzcd}[ampersand replacement=\&]
E' \arrow[r,"\beta"] \& \unsure{\mathcal{P}|\mathscr{P}}(E)
\end{tikzcd}}
% source lines 7--10
\unsure{i.e.} une application
\diagram{\begin{tikzcd}[ampersand replacement=\&]
\mathcal{P}(E') \arrow[r,no head,"\bar\beta"] \& \mathcal{P}(E)
\end{tikzcd}}
commutant aux réunions. On va supposer que
\[
\forall x'\in E',\quad \bar\alpha\bigl(\beta(x')\bigr)\subset\operatorname{supp}^{\circ}_{E'}(\{x'\}). \tag*{$(*)$}
\]
% source lines 11--16
Si $B\subset E$, \deleted{alors} et $x'\in E'$, alors
\[
x'\mid\circ\mid\bar\alpha(B)\Longleftrightarrow\bar\alpha\bigl(\beta(x')\bigr)\mid\circ\mid\bar\alpha(B)
\]
% margin note written slanted
\margin{\uline{NB} \unsure{si l'hypoth.} \unsure{$(*)$} \unsure{est une hyp.} \unsure{de fidélité}}
et, si $\bar\alpha$ a une propriété de \uline{fidélité} pour $\mid\circ\mid$, on en conclut
\[
\beta(x')\mid\circ\mid B
\]
% source lines 17--21
Appliquons ceci à la situation $(**)$ (p. \unsure{précéd.}) à $(\unsure{B},x')$ et à $(C,y')$, on trouve
\[
\beta(x')\mid\circ\mid B;\quad
\underbrace{\beta(y')\mid\circ\mid C=\operatorname{cosupp}^{\circ}(B)}_{\text{i.e.}\quad\beta(y')\subset\operatorname{supp}^{\circ}(B)}
\]
d'où
\[
\beta(x')\mid\circ\mid\beta(y'),\qquad\text{d'où}
\]
% archive sheet number: 65

\scanpage{67}
\mspage{63}
% source lines 1--6
\[
\bar\alpha(\beta(x'))\mid\circ\mid\bar\alpha(\beta(y'))
\]
(\unsure{comp.} de $\alpha$ aux disjonctions), pour en conclure, comme on voudrait, $x'\mid\circ\mid y'$, il faudrait que l'inclusion dans $(*)$ soit une identité; ou plutôt que
\[
\overset{\deleted{\unsure{E}}}{\operatorname{supp}}^{\circ}_{E'}(\bar\alpha(\beta(x'))=\operatorname{supp}^{\circ}_{E'}(\{x'\}) .
\]
% source lines 7--11
Voici un choix \uline{optimal} de $\beta$, \unsure{en termes} de $\alpha$:

Soit $\beta=\deleted{\unsure{t}}\alpha^*$ défini par
\[
\deleted{\unsure{\beta}}\alpha^*(x')=\{x\in E\mid\alpha(x)\subset\operatorname{supp}^{\circ}_{E'}(\{x'\})\}
\]
% source lines 12--16: abandoned formulations
\deleted{Notre hypothèse supplémentaire sur}
\deleted{Avec cette \unsure{définition}\inserted{ce choix de $\beta$}, on aura}
\deleted{Avec cette}
\unsure{On aura} alors tautologiquement
\[
\bar\alpha(\alpha^*(x'))\subset\operatorname{supp}^{\circ}_{E'}(\{x'\}),
\]
% source lines 17--20
\deleted{et il faut supposer que cette \unsure{inclusion|identité}}
\deleted{\unsure{de fait|du fait}}\inserted{\unsure{donc}}
\[
\operatorname{supp}^{\circ}_{E'}(\bar\alpha(\alpha^*(x')))\subset\operatorname{supp}^{\circ}_{E'}(\{x'\})
\]
et il faut supposer que cette inclusion est une \uline{égalité}
\[
\boxed{\operatorname{supp}^{\circ}_{E'}(\bar\alpha(\alpha^{\unsure{*|t}}(x')))=\operatorname{supp}^{\circ}_{E'}(\{x'\}).}\ .\tag*{(i)}
\]
% source lines 21--25
(\unsure{Comparer} l'hypothèse de surjectivité dans l'exemple de \unsure{tantôt}). De plus, la propriété que $\alpha$ soit compatible, doit être remplacée par une hyp. de ``compatibilité fidèle''
\[
\boxed{\text{Si }x,y\in E,\quad x\mid\circ\mid y\Longleftrightarrow\bar\alpha(x)\mid\circ\mid\bar\alpha(y)}\tag*{(ii)}
\]
% archive sheet number: 66

\scanpage{68}
\mspage{64}
% source lines 1--2
\unsure{qui|que} implique aussitôt, pour $A,B\subset E$
\[
\bar\alpha(A)\mid\circ\mid\bar\alpha(B)\Longleftrightarrow A\mid\circ\mid B.
\]
% source lines 3--7
\unsure{Ces} conditions draconiennes donnent ce qu'il nous faut, mais peut-être même un peu \deleted{plus} trop --- \unsure{n'impliqueraient-elles} pas qu'on a
\diagram{\begin{tikzcd}[ampersand replacement=\&]
\alpha_*:\Sigma_E \arrow[r,"\sim"] \& \Sigma_{E'}\quad\text{bijectif ?}
\end{tikzcd}}
% source lines 8--14
\deleted{\illegible{2}}\inserted{D'autre} part, on aimerait peut-être \deleted{\unsure{mais|avoir}} avoir la formule $(*)$ p. 59, dans le cas où la formule $(**)$ p. 59 est fausse (comme ds l'exemple de tantôt). Donc je n'ai pas l'impression encore d'avoir vraiment bien compris la situation.

% source lines 15--17
Dans la formule $(*)$ p. 59
\[
\underbrace{\alpha_*\operatorname{supp}^{\unsure{\circ|\bullet}}_E(A)}_{\substack{\Vert\,\text{déf}\\\operatorname{supp}^{\circ}_{E'}\bar\alpha(\operatorname{supp}^{\circ}_E(A))}}
=\operatorname{supp}^{\circ}_{E'}(\bar\alpha(\unsure{A}))
\]
% source lines 18--21
l'inclusion $\supset$ est claire, puisque
\[
\operatorname{supp}^{\circ}_E(A)\supset\deleted{\operatorname{supp}^{\circ}_E}A\quad\text{d'où}\quad\deleted{\unsure{\alpha}}
\]
\[
\bar\alpha(\operatorname{supp}^{\circ}_E(A))\supset\bar\alpha(A)
\]
et il suffit de passer aux $\operatorname{supp}^{\circ}_{E'}$. Donc le p\textsuperscript{b} est
% archive sheet number: 67

\scanpage{69}
\mspage{65}
% source lines 1--3
\[
\begin{aligned}
\alpha_*(\operatorname{supp}^{\unsure{\circ|\bullet}}_E(A))&\overset{?}{\subset}\operatorname{supp}^{\circ}_{E'}(\bar\alpha(A))\\
\Vert\qquad&\\
\operatorname{supp}^{\circ}_{E'}(\bar\alpha(\operatorname{supp}^{\circ}_E(A))&
\end{aligned}
\]
ou encore de ceci:
\[
\left\{\begin{aligned}
&\text{Si }B\subset\operatorname{supp}^{\circ}_E(A)\quad\text{i.e.}\quad\forall x\in E,\ x\mid\circ\mid A\Rightarrow x\mid\circ\mid B\\
&\text{peut-on en conclure que}\\
&?\left\{\begin{aligned}
\bar\alpha(B)&\subset\deleted{\unsure{\bar\alpha(\operatorname{supp}}}\operatorname{supp}^{\circ}_{E'}(\bar\alpha(A))\quad\text{i.e.}\quad\forall x'\in E'\\
&x'\mid\circ\mid\bar\alpha(A)\Rightarrow x'\mid\circ\mid\bar\alpha(B
\end{aligned}\right.\\
&\hfill ?
\end{aligned}\right.
\]
% source lines 8--12
Pour l'affirmer, il faudrait une corr. en sens inverse
\diagram{\begin{tikzcd}[ampersand replacement=\&]
\beta:E' \arrow[r,no head] \& \unsure{\mathcal{P}|\mathscr{P}}(E)
\end{tikzcd}}
telle que l'on ait, $\forall x\in E,\ x'\in E'$
\[
x\mid\circ\mid\beta(x')\Longleftrightarrow\alpha(x)\mid\circ\mid x',\tag*{$(*)$}
\]
\inserted{d'où, pour $A\subset E,\ B'\subset E'$
\[
A\mid\circ\mid\bar\beta(\deleted{A'}\unsure{B'|A'})\Longleftrightarrow\bar\alpha(A)\mid\circ\mid\deleted{A'}B'\tag*{$(**)$}
\]}
% source lines 13--17
Dans ce cas, l'implication ? ci-dessus se \deleted{\unsure{lit}} reformule comme
\[
\deleted{\unsure{x'|x}\mid\circ\mid\beta(\unsure{x'|x})\overset{?}{\Longrightarrow}}
\]
\[
\beta(x')\mid\circ\mid A\overset{?}{\Longrightarrow}\beta(x')\mid\circ\mid B
\]
% margin note written slanted
\margin{On n'a \unsure{pas} utilisé ici $x\mid\circ\mid y\Rightarrow\alpha(x)\unsure{\mid\circ\mid}\alpha(y)$!}
qui résulte bien de $B\subset\operatorname{supp}^{\circ}_E(A)$

% source lines 18--23
Ici encore, si une telle corr. $\beta$ existe, il doit y en avoir une optimale: soit en effet
\[
\gamma(x')\unsure{\equiv|=}\{x\in E\mid\alpha(x)\mid\circ\mid x'\},
\]
\margin{cette fois \unsure{définie|défini} en termes de $\alpha$ seul}
alors la propriété $(*)$ sur $\beta$ se lit
\[
x\mid\circ\mid\beta(x')\Longleftrightarrow x\in\gamma(x')\qquad\text{i.e.}
\]
\[
\gamma(x')=\operatorname{cosupp}^{\circ}(\beta(x'))\qquad\deleted{\text{\unsure{on aura}}}
\]
% archive sheet number: 68

\scanpage{70}
\mspage{66}
% source lines 1--4
ce qui signifie deux choses
\[
\begin{aligned}
1^\circ)&\quad\gamma(x')\in\Sigma_{\unsure{E}}\\
2^\circ)&\quad\deleted{\unsure{\beta}}\operatorname{supp}^{\circ}\beta(x')=C\gamma(x')\quad(=\operatorname{cosupp}^{\circ}(\gamma(\unsure{x'|x})))
\end{aligned}
\]
Donc une telle $\beta$ existe ssi on a $1^\circ)$, alors on peut prendre
\[
\begin{aligned}
\beta(x')&=\operatorname{cosupp}^{\circ}(\gamma(x'))=\{x\in E\mid\forall y\in E,\ \alpha(y)\mid\circ\mid x'\\
&\hspace{7em}\Longrightarrow x\mid\circ\mid y\}
\end{aligned}
\]
% margin note written slanted
\margin{\uline{NB} Si $E,E'$ sont des disjoncteurs \unsure{triviaux} ($x\mid\circ\mid y\Leftrightarrow x\ne y$) alors on voit aussitôt que $\beta$ associé à un $\alpha$ est unique (car $\operatorname{supp}^{\circ}_E=\mathrm{id}$) et n'est autre que la \unsure{correspondance} \uline{transposée} de la corr. de départ. Donc la relation entre $\alpha$ et $\beta$ envisagée ici est une \unsure{généralisation} de la notion de transposée d'une correspondance.}
% source lines 7--14
\unsure{Cette $\beta$ est la|C'est la} ``plus grande'' de toutes celles satisfaisant $(*)$ [et on l'obtient \inserted{\unsure{ainsi|aussi}}: d'une $\beta$ quelconque, en composant avec la fonction $\operatorname{supp}^{\circ}_E$].

\unsure{On note|On notera} d'ailleurs que \deleted{cette $\beta$} le $\beta$ associé à $\alpha$ ne change pas, quand on remplace $\alpha$ par son composé avec $\operatorname{supp}^{\circ}_{E'}$.

% source lines 15--21
Je vois se dégager une proposition.

% The following draft block is crossed through diagonally; it also contains earlier horizontal deletions.
\deleted{\uline{Proposition} Soient $(E,\unsure{\mid\circ\mid_E})$, $(E',\mid\circ\mid_{E'})$ deux disjoncteurs. \deleted{Il y a \unsure{donc} \illegible{3}}\inserted{Soit $\operatorname{Corr}(E,E')$ l'ensemble} \deleted{\unsure{corres}} des correspondances de $E$ dans $E'$, et Soit $\operatorname{Corrdis}(E,E')$ le sous-ensemble des}
% archive sheet number: 69

\scanpage{71}
\mspage{67}
% source lines 1--8: cancelled draft, with horizontal and diagonal strokes
\deleted{\unsure{Ce|ce} qui signifie \unsure{d}}

\deleted{\uline{Définition}, Un couple $(\alpha_*,\beta_*)$ d'applications}
\deleted{\inserted{application}}
\diagram{\begin{tikzcd}[ampersand replacement=\&]
\deleted{\alpha_*:\Sigma_E} \arrow[r] \& \deleted{\Sigma_{E'}}
\end{tikzcd}}
\deleted{est dite une \uline{correspondance} \uline{\unsure{disjonction|disjonctive}} de $E$ à $E'$, si \deleted{il}\inserted{elle} \deleted{existe une application}\inserted{\unsure{commute aux}}}
\[
\deleted{\beta_*:\Sigma_{E'}\unsure{:}}
\]
\deleted{$\operatorname{Sup}$ quelconques, et s'il existe une application}

% source lines 9--15
\uline{\unsure{Suite|Sottise}} Une correspondance $\alpha$ entre deux ensembles $E,E'$ (i.e. une \uline{relation} entre $E,E'$ i.e. une partie $\alpha\subset E\times E'$) peut être interprétée comme une application
\diagram{\begin{tikzcd}[ampersand replacement=\&]
\alpha_*:\unsure{\mathcal{P}|\mathscr{P}}(E) \arrow[r] \& \mathcal{P}(E')
\end{tikzcd}}
qui commute aux $\bigcup_i$ quelconques.
% margin note written slanted
\margin{Composition des corr. $\Leftrightarrow$ \unsure{composition} des applications \unsure{sur|les} $\mathcal{P}(\unsure{E|E'})$.}

% source lines 16--22
De même, \deleted{\unsure{une}} si $E,E'$ sont des disjoncteurs, je vais considérer les applications
\diagram{\begin{tikzcd}[ampersand replacement=\&]
\alpha_*:\Sigma_E \arrow[r,no head] \& \unsure{\Sigma}_{E'}
\end{tikzcd}}
qui commutent aux $\operatorname{Sup}$ quelconques, et qui généralisent donc les correspondances de disjoncteurs \unsure{``}triviaux''. Il suffit de se donner l'application
\diagram{\begin{tikzcd}[ampersand replacement=\&]
\alpha:E \arrow[r] \& \Sigma_{E'}\\
x \arrow[r,maps to] \& \alpha_*(\underbrace{\operatorname{supp}^{\circ}_E(\{x\})}_{\unsure{S|s}(x)})
\end{tikzcd}}
% archive sheet number: 70

\scanpage{72}
\mspage{68}
% source lines 1--5
et on \deleted{pourra ainsi} l'étend en $\alpha_*$ par
\[
\alpha_*(S)=\underset{x\in S}{\operatorname{Sup}\text{ dans }\Sigma_{E'}}
(\underbrace{\alpha(x)}_{\substack{\Vert\\\alpha_*(S_{\deleted{\unsure{E}}}(x))}})
=\operatorname{supp}^{\circ}_{E'}(\bar\alpha(S))
\]
\uline{formule valable}, puisque
\[
S=\left[\underset{x\in S}{\operatorname{Sup}_{\deleted{\unsure{E}}}(\text{dans }\Sigma_E)}\right](S_x).
\]
Ainsi

% source lines 6--11
\uline{Lemme} 1 Corr. biunivoque entre \deleted{\unsure{les}} \deleted{Appl} $\operatorname{Hom\,ens}(E,\Sigma_{E'})$ et $\operatorname{Hom\,ord}^{!}(\Sigma_E,\Sigma_{E'})$, où le $!$ signifie qu'on prend les hom. d'ens. ordonnés commutant aux $\operatorname{Sup}$ quelconques.
% margin note written slanted
\margin{Mais il n'est pas clair qu'on ait la formule, pour tt $A\subset E$ $\deleted{\unsure{\beta|\Delta}}\alpha_*(\operatorname{supp}^{\circ}_E(A))\overset{?}{=}\operatorname{supp}^{\circ}_{E'}(\bar\alpha(A))$ mais seulement que $\supset$}

% source lines 12--18
On voudrait cependant une notion de \unsure{transposé|transposée} d'un tel \unsure{fourbis}, qui généralise le cas connu des ``disjoncteurs triviaux'', où la définition du transposé $\beta$ de $\alpha$ peut se décrire par
\[
\deleted{\alpha(x)\ne}
\]
\[
x'\unsure{\in|\notin}\alpha(x)\Longleftrightarrow\unsure{\alpha|x}\in\beta(x')
\]
% source lines 19--21
ici on passe à la négation
\[
x'\mid\circ\mid\alpha(x)\Longleftrightarrow x\unsure{\mid\circ\mid}\beta(x')
\]
On a donc
% archive sheet number: 71

\scanpage{73}
\mspage{69}
% source lines 1--5
\uline{Lemme 2} Considérons une application
\diagram{\begin{tikzcd}[ampersand replacement=\&]
\alpha_*:\Sigma^{*}_{\unsure{A|E}} \arrow[r] \& \Sigma^{*}_{A'}
\end{tikzcd}}
\inserted{\deleted{\unsure{déduite}}\deleted{\unsure{par|d'}} ``\uuline{associée}'' à une correspondance \deleted{\unsure{comme de $E'$}}\unsure{,|)}}
\unsure{donne|donné|donnée} dans le lemme 1,
\diagram{\begin{tikzcd}[ampersand replacement=\&]
\alpha:E \arrow[r] \& \unsure{\mathcal{P}|\mathscr{P}}(E')
\end{tikzcd}}
en composant avec $\operatorname{supp}^{\circ}_{E'}$, donc
\deleted{équivalentes}
\[
\alpha_*(S)=\operatorname{supp}^{\circ}_{E'}(\bar\alpha(S)).
\]
Conditions équivalentes
% source lines 6--9
\diagram{\begin{tikzcd}[ampersand replacement=\&]
\text{(i)}\quad\exists\ \beta_*:\Sigma^{\unsure{\prime|*}}_{A'} \arrow[r] \& \Sigma^{\prime}_{A},\quad\text{commutant}
\end{tikzcd}}
aux $\operatorname{Sup}$ quelconques (i.e. comme $\alpha_*$) et satisfaisant, pour \deleted{tt} $S\in\Sigma^{\unsure{\wedge|*}}_A,\ S'\in\Sigma^{\wedge}_{A'}$
\[
\alpha_*(S)\mid\circ\mid S'\Longleftrightarrow S\mid\circ\mid\beta_*(S')\tag*{$(*)$}
\]
% margin note written slanted
\margin{\unsure{Autres conditions} équivalentes}
% source lines 10--13
\diagram{\begin{tikzcd}[ampersand replacement=\&]
\text{(ii)}\quad\exists\ \beta:E' \arrow[r] \& \deleted{\Sigma^{*}_E}\mathcal{P}(E),\quad\text{telle que}
\end{tikzcd}}
l'on ait, pour $x\in E,\ x'\in E'$
\[
\alpha(x)\mid\circ\mid\unsure{y'|x'}\Longleftrightarrow x\mid\circ\mid\beta(x').\tag*{$(**)$}
\]
% source lines 14--17
De plus, l'application $\beta_*$ dans (i) est uniquement déterminée, et si $\beta$ est comme dans (ii), $\beta_*$ est \uline{associé} à $\beta$
\[
\beta_*(\unsure{S'|B'})=\operatorname{supp}^{\circ}_E(\bar\beta(S'))
\]
% source lines 18--22
Les conditions (i) (ii) sont \unsure{équivalentes à la} suivante.
% margin note written slanted
\margin{Autres conditions ? équivalentes: (iv) $\forall S'\in\Sigma_{E'}$, l'ens $\{x\in X\mid\bar\alpha(x)\subset S'\}$ est $\in\Sigma_E$ (v) Validité de la formule $\alpha_*(\operatorname{supp}^{\circ}_E(A))\unsure{\overset{?}{=}|\ne}\operatorname{supp}^{\circ}_{E'}(\bar\alpha(A))$}
\[
\text{(iii)}\quad\forall x'\in E',\ \text{soit}
\]
\[
\bar\beta(x')=\{x\in E\mid\alpha(x)\mid\circ\mid x'\}\subset E
\]
\deleted{Alors la partie \unsure{\texttt{|}}}
On veut
\[
\bar\beta(x')\in\Sigma_E.
\]
% archive sheet number: 72
